
\chapter{Мультивалютность и трансграничные платежи}~\label{ch:multicurrency}
Один платёж в сервисе бронирования легко распадается на четыре
валюты: покупатель видит цену в евро, платит картой в фунтах,
мерчант получает выплаты в долларах, а расчёт с эквайером приходит
в четвёртой валюте. Каждая сумма живёт в своём временном слое.
Корреспондентские счета и трансграничная логика карточных сетей
разобраны в предыдущих частях; здесь фокус на продуктовых
решениях поверх контура: кто задаёт курс, где появляется
валютный риск, в какой валюте живут обязательства.

\section{Кто управляет курсом}~\label{sec:mc-currencies}

Первое архитектурное решение в мультивалютном продукте:
кто фиксирует курс, который увидит клиент
и на который потом будет ссылаться поддержка. Вариантов четыре.

Первый вариант: курс фактически остаётся на стороне
эмитента и правил карточной сети. Так устроен стандартный
карточный сценарий: мерчант выставляет цену в валюте операции,
а сумма в валюте карты определяется уже на стороне эмитента
по правилам валюты списания и конверсии~\cite{visa-core-rules,mc-dcc-guide}.
Для мерчанта это просто, но ни мерчант, ни платформа заранее
не знают, какой окажется итоговая сумма списания в валюте карты.

Второй вариант: курс контролирует эквайер или DCC-провайдер.
Тогда держателю карты показывают итоговую сумму в домашней
валюте ещё до авторизации. Цена становится предсказуемой
для покупателя, но право задать курс и удержать маржу
переходит к стороне, которая управляет DCC~\cite{visa-core-rules,mc-dcc-guide}.

Третий вариант: курс задаёт сама платформа. Так может быть
устроен кошелёк, сервис путешествий или маркетплейс
с собственным валютным слоем: интерфейс показывает курс,
платформа берёт на себя риск его устаревания на коротком
окне, а затем хеджирует или неттирует фактические потоки
на этапе расчёта.

Четвёртый вариант: курс приходит от внешнего казначейского
или валютного провайдера, а платёжный продукт лишь привязывает
платёж к уже рассчитанной сделке. Этот сценарий типичен
для B2B-сервисов и платформ выплат, где цена перевода живёт
дольше одной сессии оплаты.

Курс нужно моделировать как отдельное обещание
с явным владельцем, сроком действия и источником. Как только
эта роль размыта, спор о потерях на конверсии почти неизбежно
переходит из архитектуры в поддержку и сверку.

Практическая проблема -- устаревание курса. Если платформа показывает
клиенту курс в момент выбора, а capture приходит через час, курс
за это время мог сдвинуться. Разница между показанным и фактическим
курсом ложится на того, кто зафиксировал обещание. Обычное
решение: курс действует в пределах TTL (5--15 минут); если
capture приходит позже, платформа либо пересчитывает по
актуальному курсу, либо принимает риск на себя. Без явного TTL
потери на устаревшем курсе накапливаются незаметно, пока
казначейство не обнаружит их при ежемесячной сверке.

\begin{figure}[tbp]
  \centering
  \includefiguresvg[width=.8\linewidth]{assets/figures/ch29-multicurrency/fx-ownership-models}
  \caption{Четыре модели владения курсом.}~\label{fig:fx-ownership-models}
\end{figure}

Множественные транзакции усиливают эффект. Подписка с ежемесячным
списанием в EUR при выплате мерчанту в USD проходит через
конверсию 12 раз в год. Каждый раз курс другой, и
совокупная дельта может составить заметную долю маржи платформы.
Поэтому зрелые продукты фиксируют курс как отдельный объект:
источник, значение, момент фиксации, TTL и сторона, несущая риск.

\practicenote{%
Если продукт обещает клиенту \enquote{ваш курс}, у этого курса должны
быть явный источник и окно действия. Иначе потери на конверсии быстро
становятся спором между поддержкой, финансами и сверкой.

}
\section{Где фиксируется валюта расчёта}~\label{sec:mc-settlement}

В какой валюте система считает свои обязательства -- отдельный
вопрос. В нём участвуют три сущности.

\textbf{Валюта цены (transaction currency)}: валюта,
которую видит покупатель и которую мерчант ожидает
на странице оплаты.
\textbf{Валюта списания (billing currency)}: валюта,
в которой эмитент дебетует держателя карты.
\textbf{Валюта расчёта (settlement currency)}: валюта,
в которой сеть или эквайер фактически рассчитывается с вашей
стороной~\cite{visa-core-rules,mc-rules}.

Архитектурная ошибка начинается там, где валюта расчёта
считается производной от валюты цены по умолчанию.
На практике продукт обычно выбирает одну из двух моделей.

\textbf{Единая валюта расчёта.} Внутренний реестр обязательств
и выплаты ведутся в одной базовой валюте. Это упрощает
операционный контур: один казначейский баланс, одна логика
комиссий, меньше сценариев для сверки. Цена этой простоты:
постоянная конверсия на входе или выходе из системы
и зависимость от валютной наценки внешнего провайдера.

\textbf{Несколько валют расчёта.} Система хранит обязательства
и расчёты в нескольких валютах. Это сложнее для реестра
обязательств и логики выплат, но избавляет от лишней
конверсии и не смешивает реальные валютные позиции
в одну синтетическую сумму.

Вопрос при проектировании: \enquote{EUR здесь --
  валюта цены, валюта списания,
валюта расчёта, валюта хранения обязательства или всё сразу?}
Ответы часто не совпадают, и из этого расхождения
следует большая часть архитектурных компромиссов.

Конкретный пример: французский покупатель платит картой в EUR
за бронирование на британской платформе (GBP). Эквайер
рассчитывается с платформой в USD. Валюта цены -- EUR, валюта
списания -- EUR (совпадает с валютой карты), валюта расчёта -- USD.
Если платформа выплачивает мерчанту в GBP, появляется ещё одна
конверсия: USD~→~GBP на этапе выплаты. Реестр обязательств должен
хранить все три точки фиксации: сумму в EUR на момент покупки,
сумму в USD на момент settlement и сумму в GBP на момент выплаты.
Без этого сверка не сможет объяснить, почему мерчант получил не ту
сумму, которую видел покупатель.

\begin{figure}[tbp]
  \centering
  \includefiguresvg[width=.8\linewidth]{assets/figures/ch29-multicurrency/three-currencies-timeline}
  \caption{Три валюты одной операции: цена, расчёт, выплата.}~\label{fig:three-currencies-timeline}
\end{figure}

\section{DCC и маршрутизация по коридорам}~\label{sec:mc-dcc}

DCC показывает держателю карты сумму в знакомой валюте,
но одновременно меняет логику конверсии:
расчёт по сети идёт по одному маршруту, а сумма списания для клиента --
по другому~\cite{visa-core-rules,mc-tpr}.

Продукту нужно хранить исходную сумму, показанную
сумму и применённый курс как отдельные значения, а не пытаться
вывести одно из другого задним числом. Без этого невозможно
ни объяснить клиенту списание, ни сверить удержанную DCC-маржу~\cite{visa-core-rules,mc-tpr}.

DCC должен быть частью политики маршрутизации. Для одних
коридоров он может быть экономически оправдан, для других
требует особенно осторожной оценки из-за репутационного риска.
Внутренняя логика здесь обычно оценивает коридор, BIN,
страну эквайера и тип мерчанта, а затем решает, допустим ли
DCC и откуда брать курс.

DCC нельзя проектировать отдельно от комплаенс-контура.
Если согласие клиента, раскрытая наценка и след аудита
не фиксируются как часть события, позже будет трудно доказать,
что конверсия была предложена корректно~\cite{visa-core-rules,mc-tpr}.

Как только у платформы появляется больше одной страны, одного
эквайера или одного платёжного контура выплат, мультивалютность
превращается в задачу маршрутизации по коридорам: для каждой пары
\enquote{страна покупателя / страна мерчанта} нужно выбрать
весь маршрут целиком: эквайера, валютный контур и способ выплаты.

Обычно этот выбор сводит в одном решении четыре вопроса:
где дешевле авторизация и расчёт, где выше доля одобрений
для нужного диапазона BIN, в какой валюте мерчант хочет получить
выплату и на чьей стороне продукт готов держать валютный риск.
На практике выбирают одну из двух моделей выплат.

\textbf{Выплата в домашней валюте мерчанта.} Мерчант всегда получает
в домашней валюте, а платформа сама консолидирует и конвертирует
все входящие потоки. Это проще для мерчанта, но сложнее
для казначейской функции платформы.

\textbf{Выплата в валюте исходного потока.} Платформа хранит
отдельные валютные карманы и платит мерчанту в той валюте,
в которой пришёл поток. Это уменьшает число конверсий,
но требует, чтобы мерчант умел жить с несколькими валютными
балансами и с отдельным графиком выплат по каждой из них.

\begin{figure}[tbp]
  \centering
  \includefiguresvg[width=.75\linewidth]{assets/figures/ch29-multicurrency/corridor-routing}
  \caption{Четыре решения в маршрутизации по коридорам.}~\label{fig:corridor-routing}
\end{figure}

Маршрутизация по коридорам связывает в один узел стоимость
авторизации, долю одобрений, валютный риск и модель выплат.
Разрывать эти решения по разным командам опасно: финансовая
логика быстро начинает спорить с протокольной.

\practicenote{%
Если внутри принимаются потоки в пяти валютах, а мерчанту обещан один
баланс, это казначейская функция.
У валютной позиции, лимитов и источников курса должен быть явный владелец.

}
\section{Сверка по валютам}~\label{sec:mc-fx-risk}

Расхождения в мультивалютном продукте появляются
не на этапе авторизации, а позже, когда
одна и та же операция начинает жить сразу в нескольких
представлениях: на странице оплаты в валюте цены, в журнале
авторизации с одной суммой, в клиринге уже с другой
валютой расчёта, а в отчёте о выплате после удержания
комиссии и валютного спреда. Валютную сверку нужно
проектировать как отдельный слой,
а не как частный случай обычной сверки. Карточные правила
прямо требуют сверять данные сети со своими внутренними
записями; в Mastercard это сформулировано как ежедневная
обязанность сверять итоги и транзакции с внутренним учётом~\cite{mc-rules}. Три правила.

Первое: не терять исходную валютную пару. Внутренний реестр
обязательств должен хранить исходную сумму, сумму после
конверсии, источник курса и момент его фиксации.

Второе: считать комиссию отдельно от конверсии. Если курсовая разница
смешана с эквайерской комиссией, операционная команда никогда
не поймёт, где именно исчезла маржа.

Третье: сверять выплаты по валютным карманам, а не по агрегированной
сумме. Сходящаяся общая цифра легко скрывает дыру внутри отдельной
валюты, особенно если одна валюта перекрывает другую по номиналу.

Пример валютной сверки одной операции. Покупка за 100~EUR,
settlement по курсу 1,08 = 108~USD, комиссия эквайера
2,5\% = 2,70~USD, interchange 1,5\% = 1,62~USD, комиссия
платформы 3\% = 3,24~USD, выплата мерчанту =
108 $-$ 2,70 $-$ 1,62 $-$ 3,24 = 100,44~USD. Если мерчант
получает выплату в GBP по курсу 0,79, итого 79,35~GBP. Одна
и та же покупка существует в пяти представлениях: 100~EUR на
странице оплаты, 100~EUR в авторизации, 108~USD в settlement,
100,44~USD в нетто-выплате, 79,35~GBP на счёте мерчанта.
Каждое представление корректно, но ни одно не равно другому.

Ошибки сверки чаще всего возникают при смешении валютной
дельты и комиссий. Если 7,70~USD удержаний
(2,70 + 1,62 + 3,24) учтены как FX loss, а не как
раздельные комиссии, казначейство видит завышенную курсовую
дельту и начинает искать проблему там, где её нет. Каждое
удержание должно иметь свой тип: interchange, assessment,
platform fee, FX markup -- и сверяться по типу, а не по
итоговой сумме.

\begin{figure}[tbp]
  \centering
  \includefiguresvg[width=.85\linewidth]{assets/figures/ch29-multicurrency/multicurrency-recon}
  \caption{Одна транзакция в пяти валютных представлениях.}~\label{fig:multicurrency-recon}
\end{figure}

Мультивалютность почти всегда тянет за собой
изменения в реестре обязательств из главы~\ref{ch:internal-ledger}.
Без отдельных балансов, модели начисления комиссий и связи
с графиком выплат валютная сверка быстро становится ручной
и хрупкой.

\section{Корреспондентский банкинг и ISO~20022}~\label{sec:mc-correspondent}

Корреспондентские расчёты разобраны в главе~\ref{ch:banks}, миграция к ISO~20022 -- в главе~\ref{ch:open-banking}. Для архитектурных решений здесь достаточно
трёх инвариантов.

Первый: расчёт по трансграничной операции почти всегда проходит
по более длинной цепочке, чем авторизация, и живёт
в другом тайминге~\cite{fatf-correspondent-banking,swift-instant-interoperability}. Одобрение за секунду не означает выплату
в тот же день.

Второй: расчётный банк, карточная сеть и валютный провайдер
могут быть тремя разными сторонами. Продукту нужно
явное отображение того, кто отвечает за каждую стадию,
а не абстрактное \enquote{деньги идут через SWIFT}.

Третий: переход отрасли на ISO~20022 делает сообщения богаче,
но не снимает продуктовых решений. Структурированные поля
облегчают комплаенс-проверки и сверку, однако сами по себе
не решают, кто несёт валютный риск и в какой валюте должны
жить внутренние обязательства~\cite{iso-20022,swift-iso20022-milestone,swift-transaction-manager-iso20022}.

Инфраструктура делает поток наблюдаемым, но не принимает
за продукт решений о риске, курсе и модели расчёта.

